\subsection{Herbie: Making an \eqsat Application Sound}
\label{sec:herbie}


Herbie~\cite{herbie} is a widely-used, open-source tool
 for making floating-point programs more accurate,
 with thousands of users and yearly stable releases.
Herbie takes as input a real expression,
 and returns the most accurate floating-point implementation
 it can synthesize.
Since floating-point error is a critical issue
 in scientific computing,
 Herbie is used in a variety of domains,
 including machine learning, computer graphics,
 and computational biology.


The core of Herbie's algorithm
 is to run equality saturation to explore equivalent programs.
These programs are mathematically equivalent
 over the real numbers,
 but may have different behavior over floating-point numbers.
Herbie considers candidate programs from the
 results of equality saturation,
 finding the most accurate among them.

Herbie's rewrite rules are known to be unsound,
 which has been the cause of numerous bugs in the past.
In addition, unsound rules prevent Herbie from running equality saturation
  longer once unsoundness occurs.
Unfortunately, merely removing the unsound rules makes Herbie
  useless on a large portion of its benchmark suite.
For example, \autoref{fig:herbierule} shows a rewrite rule which
  is critical to Herbie's ability to find more accurate programs.
Using these unsound rules in a sound way requires a more sophisticated
  analysis of Herbie's input programs.
This analysis was nearly impossible with Herbie's existing \egraph implementation.

\Egglog has allowed us to implement precise analyses to
 make all rewrites sound.
First, we implement an interval analysis in \egglog, allowing rules to
  utilize information about the range of terms in the program (\autoref{fig:egglogsqrt}).
This unlocks a range of crucial rules involving division, including the rule shown in \autoref{fig:herbierule2}.


\begin{figure}
  \begin{minipage}{0.40\linewidth}
  \begin{subfigure}{\linewidth}
    \[
      \frac{a*b}{c} \implies \frac{a}{\frac{c}{b}}
    \]
    \caption{
      A fraction rule which requires $b \neq 0$.
    }
    \label{fig:herbierule2}
  \end{subfigure}
  \\[1em]
  \begin{subfigure}{\linewidth}
    \[
      x - y \implies \frac{x^3-y^3}{x^2+xy+y^2}
    \]
    \caption{
      A more complex rule derived from
       the factorization of $x^3-y^3 = (x-y)(x^2+xy+y^2)$.
      This is sound if either $x \neq 0 $ or $y \neq 0$.
    }
    \label{fig:herbierule}
  \end{subfigure}
  \caption{
    Herbie~\cite{herbie} uses
    rewrite rules to create program variants
    with less floating-point error from phenomena like cancellation.
    Some rules are only sound under certain conditions.
  }
  \end{minipage}
  \hfill
  \begin{minipage}{0.55\linewidth}
  \begin{lstlisting}[language=egglog]
(function lo (Math) Rational :merge (max old new))
(function hi (Math) Rational :merge (min old new))

(rule ((= e (Sqrt a)))
      ((set (lo e) (rational 0 0))))

(rule ((= e (Sqrt a))
       (= loa (lo a)))
      ((set (lo e) (sqrt loa))))
  \end{lstlisting}
\caption{
  A few example rules for interval analysis of sqrt in \egglog.
  The $lo$ relation tracks the lower bound for each term,
   and we merge lower bounds by taking the max.
  Similarly, the $hi$ relation tracks the upper bound.
  First, we know that root of anything is non-negative.
  Next, since taking the root is monotonic,
   we can propagate the bounds from the arguments
   directly to the bounds of the result.
}
\label{fig:egglogsqrt}
  \end{minipage}
\end{figure}

While the interval analysis enables many of Herbie's rules, it is not
  sufficient for some more difficult cases.
We additionally implemented a ``not equals to'' analysis,
  which leverages both the interval analysis
  as well as facts inferred during rewriting.
%This composition of two analyses is particularly difficult to implement
%  in a traditional \egraph implementation.
\egglog's support for multiple interacting analyses
  enables cleanly separating interval and $\neq$ rules;
  in other \eqsat frameworks they would be implemented
  as a fused, monolithic analysis requiring much
  significantly more complicated rules.
%\rw{Can you explain how exactly \egglog makes it easier?
%Do the analyses interact with each other via rewrites? }
The not-equals analysis allows Herbie to soundly solve one of its
  classic cancellation benchmarks: $\sqrt[3]{v+1} - \sqrt[3]{v}$.
First, the interval analysis proves that $v+1 \neq v$.
%\rw{How does it prove that? It's not obvious the intervals are disjoint.}
Next, the rule $a \neq b \implies \sqrt{a} \neq \sqrt{b}$
 implies $\sqrt[3]{v+1} \neq \sqrt[3]{v}$.
This allows us to finally apply
  the rewrite from \autoref{fig:herbierule},
  substituting $\sqrt[3]{v+1}$ for $x$ and $\sqrt[3]{v}$ for $y$,
  reducing the error of the expression from extremely high to near zero.


\begin{figure}
    \includegraphics[scale=0.5]{egglogdatasmall/errorhist.pdf}

  \caption{Graph showing the difference in error between Herbie using
    \egglog's sound analysis and Herbie using the unsound ruleset across all of Herbie's benchmark suite.
  The horizontal axis is the difference 
   in the average bits of error using 
   Herbie's unsound rules vs.\ \egglog's sound rules.
  The vertical axis is the number of benchmarks. 
  Negative values
    represent benchmarks in which Herbie found a more accurate program using \egglog's analysis.}
  \label{fig:herbieeval}
\end{figure}

\begin{figure}
  \includegraphics[scale=0.5]{egglogdatasmall/timehistsecs.pdf}
  \caption{Graph showing the difference in runtime (in seconds)  between Herbie using \egglog's sound analysis and Herbie using the unsound ruleset across all of Herbie's benchmark suite.
  The horizontal axis is the difference in time to execute the benchmark.
  The vertical axis is the number of benchmarks.
  Negative values represent benchmarks in which Herbie was faster using \egglog's analysis.}
  \label{fig:herbietime}
\end{figure}

\autoref{fig:herbieeval} shows the results
 of our evaluation of Herbie using \egglog.
Herbie has a benchmark suite of \numtests floating-point programs,
 collected from a variety of domains.
We ran Herbie on each of these programs
 using both the unsound ruleset and \egglog's sound analysis.
Using \egglog's analysis makes Herbie faster overall
  (\timeegglogminutes minutes vs \timevanillaminutes minutes).
This is because \egglog generates no unsound programs, which slow down
  Herbie's search.

In \numbettererror cases,
 Herbie using a sound analysis is actually able
 to find a more accurate program than Herbie
 using the unsound ruleset.
In \numworseerror cases,
 Herbie's unsound ruleset finds more accurate results
 than \egglog's sound analysis.
%\rw{Give intuition for the accuracy differences?}
%  ZT -- too hard, need to account for taylor and regimes etc.
There are several outliers in \autoref{fig:herbieeval}.
The point on the far left represents a benchmark
  which Herbie using the unsound ruleset is unable to solve.
This input program is $9x^4 - y^2(y^2-2)$,
  and the solution involves an algebraic rearrangement
  and fma (fused multiply-add) operation.
The point on the far right represents a program that
  overflows \egglog's rational type, which can be easily fixed in the future.
